\documentclass[12pt,a4paper]{article}
\usepackage{cmap}
\usepackage[utf8]{inputenc}
\usepackage[T1,T2A]{fontenc}
\usepackage[english,russian]{babel}
\usepackage{array,booktabs,tabularx,amsmath}
\usepackage{caption,float,graphicx,qrcode}
\usepackage{geometry,hyperref,indentfirst}
\geometry{margin=25mm}
\hypersetup{unicode=true,colorlinks=true,urlcolor=blue,linkcolor=black}
\setlength{\parindent}{1.25cm}
\setlength{\parskip}{0pt}
\setlength{\emergencystretch}{2em}
\renewcommand{\labelitemi}{\textendash}
\DeclareCaptionLabelSeparator{emdash}{\space---\space}
\captionsetup[figure]{justification=centering,singlelinecheck=false,labelsep=emdash}
\addto\captionsrussian{\renewcommand{\figurename}{Рис.}}
\newcommand{\IndexDate}{07.11.2026}
\newcommand{\IndexDocuments}{478\,833}
\newcommand{\IndexTokens}{158\,544\,765}
\newcommand{\IndexSurfaceTerms}{1\,451\,704}
\newcommand{\IndexLemmaTerms}{1\,093\,912}
\newcommand{\IndexSurfacePostings}{90\,681\,947}
\newcommand{\IndexLemmaPostings}{91\,302\,977}
\newcommand{\IndexBuffer}{64}
\newcommand{\IndexDictionaryMB}{166,54}
\newcommand{\IndexPeakMiB}{208,30}
\newcommand{\IndexSeconds}{688,26}
\newcommand{\IndexVerifySeconds}{21,39}
\newcommand{\IndexVerifyMiB}{4,43}
\newcommand{\IndexMB}{1\,718,22}
\newcommand{\IndexRuns}{174}
\newcommand{\IndexMergePasses}{2}
\newcommand{\IndexSizeRows}{%
terms.bin & 101,82 \\
postings.bin & 1\,455,88 \\
documents.bin & 30,65 \\
strings.bin & 129,87 \\
}
\newcommand{\IndexExampleRows}{%
doom & Словоформы & 7\,469 & [72, 133, 165, 595, 747, \ldots] \\
minecraft & Словоформы & 5\,824 & [272, 316, 368, 530, 536, \ldots] \\
skyrim & Словоформы & 4\,635 & [446, 665, 1118, 1310, 1419, \ldots] \\
}

\begin{document}
\hypersetup{pageanchor=false}
\begin{titlepage}
\begin{center}
{\large МОСКОВСКИЙ АВИАЦИОННЫЙ ИНСТИТУТ\\
(НАЦИОНАЛЬНЫЙ ИССЛЕДОВАТЕЛЬСКИЙ УНИВЕРСИТЕТ)}
\vspace{18pt}
{\large Факультет информационных технологий и прикладной математики}

{\large Кафедра вычислительной математики и~программирования}
\vfill
{\large\textbf{Лабораторная работа №7 по курсу\\ «Информационный поиск»}}
\vspace{12mm}

{\large\textbf{Булев индекс}}
\vfill
\begin{flushright}
\begin{tabular}{@{}r@{:\hspace{0.6em}}l@{}}
\textit{Студент} & К.\,А. Арусланов\\
\textit{Группа} & М8О-403Б-23\\
\textit{Преподаватель} & А.\,А. Кухтичев\\
\textit{Дата} & \IndexDate\\
\textit{Оценка} &
\end{tabular}
\end{flushright}
\vfill
Москва, 2026
\end{center}
\end{titlepage}
\hypersetup{pageanchor=true}
\setcounter{page}{2}
\tableofcontents
\newpage

\section{Постановка задачи}
Построить булев индекс корпуса документов для последующего поиска.
Индекс должен сопоставлять термину список документов, в которых он
встречается, исключать повторные ID и обеспечивать доступ к спискам
без последовательного просмотра исходных текстов.
Структуры индекса реализовать самостоятельно на C++ с учётом
обработки корпуса, не помещающегося целиком в оперативную память.

\section{Метод решения}
\subsection{Представление данных}
Используется фиксированный снимок русскоязычных и англоязычных публикаций
о компьютерных играх из восьми источников. Индексируются заголовки
и основной текст. Документы имеют уникальные возрастающие 64-битные ID.

Для термина $t$ хранится список
\[
L(t)=\{d:\ t\text{ встречается в документе }d\},\qquad
df(t)=|L(t)|.
\]
Список упорядочен по возрастанию ID и не содержит повторений.
Несколько вхождений слова в одной статье и совпадение в разных полях
дают одну запись. Поэтому $df$ отличается от частоты вхождений,
рассчитанной при изучении закона Ципфа.

Индекс содержит словарь терминов, последовательные списки документов
и таблицу метаданных. В записи словаря сохраняются смещение списка
и число его элементов. Заголовки, URL и источники доступны по ID
документа без обращения к MongoDB.

\subsection{Построение с ограничением памяти}
Программа написана на C++17 без готовых контейнеров STL для структур
индекса. Из текста формируется поток пар «термин, ID документа».
Собственная хеш-таблица исключает одинаковые пары в текущем буфере.
При заполнении буфера пары сортируются собственной пирамидальной
сортировкой и сохраняются в промежуточную часть.

Части сливаются с помощью собственной двоичной кучи. При совпадении
пар записывается только одна из них. Если частей много, слияние
выполняется в несколько проходов; за проход открыто не более 24 входов.
Итоговый поток упорядочен сначала по термину, затем по ID документа.
Из него последовательно записываются словарь и списки.

Размер буфера пар составляет \IndexBuffer{} МиБ. Корпус и все списки
одновременно в память не загружаются. Дополнительная память используется
для текущего текстового поля, буферов слияния и морфологического словаря.

\subsection{Словоформы и леммы}
В одном индексе предусмотрены два независимых пространства терминов.
Первое хранит исходные нормализованные словоформы. Второе строится
с использованием словарного лемматизатора из ЛР4.
Неизвестные слова сохраняются, а при неоднозначном разборе документ
включается в списки всех различных лемм. В каждом отдельном списке
его ID остаётся уникальным.

Например, слово «стали» может добавить документ к спискам «сталь»
и «стать». Это расширение множества кандидатов; контекстное снятие
омонимии не выполняется. Исходное пространство позволяет позднее
сравнивать точные совпадения и поиск по леммам.
В отличие от частотного эксперимента здесь фиксируется наличие
термина в документе, поэтому альтернативные разборы не интерпретируются
как дополнительные исходные вхождения.

\subsection{Чтение и проверка}
Записи словаря имеют фиксированный размер и расположены
в лексикографическом порядке. Поиск выполняется бинарно на диске.
После нахождения записи список читается потоковым курсором.
Курсоры имеют независимые позиции, что позволяет в следующей работе
выполнять пересечение, объединение и разность списков.

Сохраняются версия формата, правила нормализации и контрольные суммы.
Проверка готового индекса последовательно читает данные, сверяет
целостность файлов, порядок словаря, уникальность ID и их принадлежность
множеству документов. Незавершённый индекс не допускается к чтению.

\subsection{План тестирования}
GoogleTest проверяет точные списки на контрольных корпусах, устранение
дублей между полями и промежуточными частями, работу с произвольными
64-битными ID, пустые документы и документы только с заголовком.
Проверяются пространства словоформ и лемм, неоднозначные разборы,
независимость курсоров и получение метаданных без исходного корпуса.

Синтетический корпус вызывает несколько проходов слияния.
Итоговые файлы побайтово сравниваются с индексом, построенным при
большем буфере. Проверяются также повреждения данных, несогласованный
каталог URL, незавершённое построение и защита существующего результата.

Успешно пройдены 10 GoogleTest и одна проверка запуска CLI.
Дополнительно выполнены 3 Python-теста подготовки каталога URL.
Python не участвует в построении словаря или списков документов.
Журнал приложен в \path{lab7/tests.txt}.

\section{Результаты}
\subsection{Объём индекса}
Индекс построен на \IndexDocuments{} документах,
содержащих \IndexTokens{} исходных токенов. Число токенов и размер
словаря словоформ совпали с результатами ЛР5 на том же снимке.

\begin{table}[H]
\centering
\begin{tabularx}{\linewidth}{@{}Xrr@{}}
\toprule
Показатель & Словоформы & Леммы\\
\midrule
Термины & \IndexSurfaceTerms & \IndexLemmaTerms\\
Пары «термин, документ» & \IndexSurfacePostings & \IndexLemmaPostings\\
\bottomrule
\end{tabularx}
\end{table}

Словарь лемм меньше словаря словоформ, однако число пар «лемма, документ»
несколько больше. Это связано с включением документа в списки всех
различных разборов: уменьшение числа терминов не гарантирует уменьшения
суммарного числа элементов списков.

\begin{table}[H]
\centering
\begin{tabularx}{\linewidth}{@{}Xr@{}}
\toprule
Файл & Размер, МБ\\
\midrule
\IndexSizeRows
\midrule
Всего & \IndexMB\\
\bottomrule
\end{tabularx}
\end{table}

ID в списках хранятся без сжатия, по 8 байт. В общий размер входят
оба пространства терминов и метаданные документов. Поэтому этот размер
не следует сравнивать с индексом только словоформ или со сжатым представлением.
Сжатие списков и ускорение их пересечения относятся к последующим работам.

\subsection{Примеры списков документов}
Программа позволяет обратиться к отдельному термину без чтения корпуса.
В таблице приведены реальные документные частоты и первые пять ID
из сохранённых списков пространства словоформ. Многоточием обозначено
продолжение списка; ID относятся к исследуемому снимку.

\begin{table}[H]
\centering
\small
\begin{tabularx}{\linewidth}{@{}llrX@{}}
\toprule
Термин & Пространство & $df$ & Начало списка\\
\midrule
\IndexExampleRows
\bottomrule
\end{tabularx}
\end{table}

\subsection{Время и память}
\begin{table}[H]
\centering
\begin{tabularx}{\linewidth}{@{}Xr@{}}
\toprule
Показатель & Значение\\
\midrule
Построение C++ с чтением и записью, с & \IndexSeconds\\
Полная проверка готового индекса, с & \IndexVerifySeconds\\
Пиковый рабочий набор при построении, МиБ & \IndexPeakMiB\\
Пиковый рабочий набор при проверке, МиБ & \IndexVerifyMiB\\
Буфер пар, МиБ & \IndexBuffer\\
Морфологический словарь, МБ & \IndexDictionaryMB\\
Промежуточные части & \IndexRuns\\
Проходы слияния & \IndexMergePasses\\
\bottomrule
\end{tabularx}
\end{table}

Измерение выполнено для нативной Windows Release-сборки GCC 13.1.0.
Приведён один прогон, а не медиана нескольких измерений.
Подготовка каталога URL и вычисление SHA-256 обвязкой не включены
во время построения C++. Рабочий набор означает резидентную память
процесса; размер выделенных буферов и занятые страницы могут различаться.

Результат подтверждает возможность построения индекса выбранного снимка
с ограниченным буфером. Производительность на миллионном корпусе
этим прогоном не проверялась. При увеличении данных дополнительно
потребуются временное место на диске и проходы слияния.

\subsection{Повторение работы}
Инструкция сборки, запуска и инспекции отдельных списков приведена
в \path{doc/lab7_run.md}. Реализация сохраняет индекс и предоставляет
интерфейс его чтения. Разбор булевых выражений и поисковая выдача
реализуются в ЛР8; оценка релевантности не относится к проверке
целостности построенного индекса.

Исходный код:
\href{https://github.com/kirbasss/MAI-IR/tree/lab7}{github.com/kirbasss/MAI-IR/tree/lab7}.
\begin{figure}[H]
\centering
\begingroup
\hypersetup{urlcolor=black}
\qrcode[height=28mm]{https://github.com/kirbasss/MAI-IR/tree/lab7}
\endgroup
\caption{Ссылка на ветку лабораторной работы №7}
\end{figure}

\section{Вывод}
Построен булев инвертированный индекс игрового корпуса.
Повторные вхождения объединены в уникальные пары «термин, документ»,
списки упорядочены и проверены на согласованность с документным множеством.
Раздельное хранение словоформ и лемм позволяет использовать одну структуру
для сравнения режимов поиска.

Внешняя сортировка и слияние ограничивают объём рабочего буфера,
а дисковый словарь и потоковые курсоры позволяют читать индекс
без загрузки всех терминов и списков в память. Следующий этап ---
булев поиск по сохранённому индексу. Его ограничения на текущем этапе:
отсутствие сжатия и координат слов, а также возможные лишние кандидаты
при неоднозначной лемматизации.
\end{document}
